\documentclass[10pt,a4paper]{amsart}
\usepackage[margin=19mm]{geometry}
\usepackage{amsmath,amssymb,mathtools}
\usepackage[scr=rsfs,cal=euler]{mathalfa}
\usepackage{xfrac}
\usepackage[unicode,hidelinks,bookmarks=false]{hyperref}
\newcommand{\ie}{i.e.\ }
\newcommand{\cf}{cf.\ }
\newcommand{\resp}{respectively\ }
\newcommand{\Subsection}{\subsection}
\providecommand{\nobreakdash}{\nobreak\hbox{-}}
\setlength{\emergencystretch}{2em}
\multlinegap0pt
\usepackage{textcomp,lmodern}
\usepackage{mathtools}
\usepackage{todonotes}
\newcommand{\added}[1]{\todo[color=green]{Added: #1}}
\newtheorem{lemma}{Lemma}[section]
\newtheorem{conjecture}[lemma]{Conjecture}
\newtheorem{corollary}[lemma]{Corollary}
\newtheorem{proposition}[lemma]{Proposition}
\newtheorem{claim}[lemma]{Claim}
\newtheorem{theorem}[lemma]{Theorem}
\newtheorem{definition}[lemma]{Definition}
\newtheorem{notation}[lemma]{Notation}
\newtheorem{question}[lemma]{Question}
\newtheorem{remark}[lemma]{Remark}
\newtheorem{example}[lemma]{Example}
\newtheorem{fact}[lemma]{Fact}
\DeclareMathOperator\DL{DL} %Diestel-Leader graph
\DeclareMathOperator\Rauz{Rauz} %Rauzy (di)graphs
\DeclareMathOperator{\rankg}{rk}% rank of a vertex
\DeclareMathOperator{\der}{der}% derangement of a path
\DeclareMathOperator{\supp}{supp}% support (of a measure)
\DeclareMathOperator{\pref}{pref} %prefix
\DeclareMathOperator{\suff}{suff} %suffix
\DeclareMathOperator{\Alg}{Alg} %von Neumann algebra
\DeclareMathOperator{\Tr}{Tr} %trace
\DeclareMathOperator{\esssup}{ess\,sup} %essential supremum 
\DeclareMathOperator{\I}{{\mathbb{I}}}
\DeclareMathOperator{\mbbP}{{\mathbb{P}}}
\DeclareMathOperator{\E}{{\mathbb{E}}}
\DeclareMathOperator{\U}{{\mathcal{U}}}
\DeclareMathOperator{\Gr}{{\mathcal{G}_*}}
\DeclareMathOperator{\Grr}{{\mathcal{G}_{**}}}
\DeclareMathOperator{\Hs}{{\mathcal{H}}}
\newcommand*\A{{\mathcal A}} %alphabet
\newcommand*\F{F} %forbidden subwords
\newcommand*\bus{{\mathfrak h}} %Busemann rank function
\newcommand*{\hprod}{\mathbin{\textnormal{\textcircled{\raisebox{0.1 em}{\tiny h}}}}}% horospheric product
\newcommand*\Lang{{\mathcal L}} %language
\newcommand*{\defi}[1]{\emph{#1}} %definition
\newcommand*{\field}[1]{\mathbf{#1}} %field
\newcommand*\orR{{\vec\Rauz}} %oriented Rauzy graph
\newcommand*{\setst}[2]{\{#1\ |\ #2\}}% set 1, such that 2
\newcommand*\Z{\field{Z}} %relative integers
\newcommand*\N{\field{N}} %integers
\newcommand*\R{\field{R}} %reals
\newcommand*\C{\field{C}} %complex 
\newcommand*\Un{\mathcal{U}} %unimodular measure
\newcommand*\eps{\varepsilon}
\begin{document}
\title[Localization of eigenfunctions in amenable unimodular random]{Localization of eigenfunctions in amenable unimodular random networks}
\author{Georgii Veprev ạte e̱gin document e̱gin abstract For an amenable unimodular random rooted network, we show that the presence of a positive point mass in the expected spectral measure implies that, with positive probability, there exists an eigenfunction with finite support. abstract; Georgii Veprev}
\date{}
\maketitle
\noindent\emph{Source and licence.} Localization of eigenfunctions in amenable unimodular random networks. Version arXiv:2606.17187v1 (CC BY 4.0). This material is distributed under the Creative Commons Attribution 4.0 International licence, http://creativecommons.org/licenses/by/4.0/, as recorded in the arXiv OAI licence metadata for this exact version and shown on the abstract page https://arxiv.org/abs/2606.17187v1. Full rights notice: licenses/arxiv-2606.17187-CC-BY-4.0.txt.\par\smallskip
\noindent\textit{Editorial scope.} Counted unit: the original \texttt{theorem} environment \texttt{theorem:balanced} at source lines 524--529 together with its complete proof at lines 530--574; the delivered document keeps source lines 267--575 so that every referenced label and every citation resolves inside it. Native editable source is the paper's own concatenated TeX; the AMS wrapper is editorial formatting only. No publisher class, upstream script or unrelated artwork is redistributed.
\par\smallskip\noindent\textit{Reference note.} This article ships no compiled bibliography (biblatex); the entries delivered here are composed from the author's own BibTeX (.bib) fields, with no field invented and no entry dropped.
\subsubsection{Amenable random networks}\label{subsection_amen}

There exist several natural notions of amenability for unimodular random rooted graphs and networks (see, e.g.,~\cite{AL07, Ka1997amenability}). Most of these conditions are equivalent under the assumption of finite expected degree. The one that is the most convenient for us in the context of this paper is the one proposed in~\cite{AL07} as it deals with a kind of F\o lner condition. 

For a subset $A\subset G $, we denote by $\partial A$ its inner vertex boundary, that is, the set of points $x \in A$ that are adjacent to a vertex of~$G$ outside~$A$. Similarly, for a non-negative integer number~$r$ we will denote by $\partial_r A$ the set of points $x \in A$ that are at (edge) distance at most~$r$ from~$G \setminus A$. In particular, $\partial_0 A = \emptyset$ and $\partial_1 A = \partial A$. We will also consider the \emph{edge-boundary $\partial_E A$}, that is, the set of edges which have one end in~$A$ and the other outside~$A$. 
%One caveat is that the F\o lner sets are, in a sense, external to the realisations of the random graph and come from different percolation processes on them. 
A \emph{percolation} on~$\rho \in \U$ is a unimodular random rooted network~$\rho^\prime \in \U(\Xi \times\Xi)$ together with a subset $\Xi_0 \subset\Xi$ such that the projection~$\pi \colon \U(\Xi \times\Xi) \to \U(\Xi)$ that forgets the second component of the marks maps~$\rho^\prime$ to~$\rho$. One can, of course, identify~$\Xi \times\Xi$ with a subset of~$\Xi$ and consider~$\rho^\prime$ as an element of~$\U(\Xi)$.  The~$\Xi_0$-component of a vertex~$x$ is the set of vertices reachable from~$x$ by edges whose second marks belong to~$\Xi_0$. We will denote by~$K(\Xi_0)$ the~$\Xi_0$-component of the root. Let~$FC(\rho)$ be the set of percolations on~$\rho$ that have only finite components almost surely. The \emph{isoperimetric constant} of~$\rho$ is the following value 
\begin{equation}\label{eq_amenability}
\imath_E(\rho) = \inf \int\frac{|\partial_E K(\Xi_0)|}{|K(\Xi_0)|} d \rho^\prime(G,o), 
\end{equation}

where $(\rho^\prime, \Xi_0)$ ranges over all percolations in $FC(\rho)$. Unimodular random rooted network~$\rho$ is called \emph{amenable} if~$\imath = 0$. In other words, $\rho$~is amenable if there exists a sequence of percolations whose clusters of the origin yield a F\o lner sequence.
%Note that the similar computation as in Proposition~\ref{proposition:rFolnler} gives for amenable~$\rho \in \U$ with finite expected degree that the same infinitum as in formula~\eqref{eq_amenability} is zero if the edge-boundary $\partial_E$ is replaced with the $r$-boundary for any positive integer~$r$.
One can also define the isoperimetric constant~$\imath_r$ for a positive integer~$r$ by replacing the edge-boundary~$\partial_E$ in formula~\eqref{eq_amenability} with the $r$-boundary~$\partial_r$. For a unimodular~$\rho$ with finite expected degree, the amenability condition $\imath_E = 0$ is equivalent to $\imath_r = 0$ for any positive integer~$r$ (see the proof of Proposition~\ref{proposition:rFolnler} in the following section).

\subsubsection{Balanced F\o lner sequences}

The definition of an amenable unimodular random network above has a disadvantage that the F\o lner sets that appear there are, in a sense, external to the realisations of the random graphs and come from different percolation processes on them. In this paper, we will use a different approach focusing on the F\o lner sequences inside the rooted networks and their properties.
Suppose that $\rho \in \mathcal{U}$ is such that for every positive integer~$n$ there exists a measurable assignment~$\{F_n(G,o)\}$ of finite subsets of~$V(G)$ such that 
\begin{equation}\label{eq:Folner}
    \lim_{n \to \infty} \int \frac{|\partial F_n (G,o)|}{|F_n(G,o)|} d\rho= 0.
\end{equation}
Naturally, we call such a sequence a \emph{F\o lner sequence} for~$\rho$. Note that the existence of a F\o lner sequence implies leaf-wise ($\rho$-almost surely) amenability of the graph. 
Different choice of the type of the boundary leads to several other versions of the F\o lner condition. We call a sequence~$\{F_n\}$ \emph{$r$-F\o lner} if the convergence~\eqref{eq:Folner} holds when~$\partial$ is replaced by~$\partial_r$ and \emph{edge-F\o lner} if~$\partial$ is replaced by~$\partial_E$. 
For graphs of uniformly bounded degree, this clearly does not give anything new. However, this does make a difference for locally finite graphs without uniform bounds, which we would like to keep under consideration.

For a measurable assignment of a subset $A(G,o) \subset G$, the function~$g_A$ is defined as
\begin{equation}\label{definition:g_function}
    g_A(G,o) = \sum\limits_{o \in A(G, x)} \frac{1}{|A(G,x)|}.
\end{equation}
The use of this function is as follows. Suppose~$\phi$ is a non-negative measurable function on $\mathcal{G_*}$. Then the average value of~$\phi$ over~$A$ can be rewritten via the mass transport principle 
\begin{multline}\label{eq:avfunction}
    \int \frac{1}{|A(G,o)|} \sum_{x \in A(G,o)} \phi(G,x) d\rho(G,o) = \\ \int \phi(G,o) \sum_{o \in A(G,x)} \frac{1}{|A(G,x)|} d\rho(G,o) = \int \phi g_A d\rho.
\end{multline}
In particular, we have $\int g_A d \rho= 1$. Let us also note that $\esssup g_A$ is the norm of the operator acting on $L^1(\Gr, \rho)$ by taking the average value over~$A$. We would like to have certain control over averages over the F\o lner sets in our graphs. As formula~\eqref{eq:avfunction} shows,  in order to do that one should impose regularity conditions on the function~$g$. This motivates the following definition.
\begin{definition}
    %Let~$\rho$ be a unimodular random network that is leaf-wise amenable. 
    We call a unimodular random network~$\rho$ ($r$-)\emph{{F\o lner} balanced} if there exists a measurable ($r$-)F\o lner sequence $F_n(G,o)$ such that for $\rho$-almost every~$(G,o)$
    $$
    % \lim_{n \to \infty} \sum\limits_{o \in F_n(G, x)} \frac{1}{|F_n(G,x)|}  = 1. 
    \lim_{n \to \infty} g_{F_n}(G,o)  = 1. 
    $$
    We call $\rho$ \emph{weakly} ($r$-)\emph{{F\o lner} balanced} if there exists a measurable ($r$-)F\o lner sequence $F_n(G,o)$ and constant $c > 0$ such that for $\rho$-almost every $(G,o)$ 
    $$
     \liminf_{n\to \infty}g_{F_n}(G,o) > c. 
    $$  
     Similarly, we define~\emph{edge-{F\o lner} balanced} and \emph{weakly edge-{F\o lner} balanced} unimodular random networks.
\end{definition}
   


The following simple proposition allows one to reduce the $r$-boundary to the standard boundary of a F\o lner set under certain regularity conditions. 
\begin{proposition}\label{proposition:rFolnler}%\todo{It might actually work without boundedness }
    Suppose that $\rho \in \mathcal{U}$ has a finite expected degree of the root. Let $\{F_n\}$ be a F\o lner sequence such that $g_{F_n}< C$ almost surely for all~$n$. Then for any positive integer~$r$
    $$
    \lim_{n \to \infty} \int \frac{|\partial_r F_n (G,o)|}{|F_n(G,o)|} d\rho= 0.
    $$
\end{proposition}
\begin{proof}
    Let us first prove the claim for~$r = 2$. Clearly $|\partial_2 A| \le \sum_{x \in \partial A} \deg x$. We have 
    \begin{multline}\label{eq:50506}
         \int \frac{|\partial_2 F_n (G,o)|}{|F_n(G,o)|} d\rho \le 
         \int \frac{\sum_{x \in \partial F_n(G,o)} \deg x}{|F_n(G,o)|} d\rho = \\
         \int \deg o\sum_{o \in \partial F_n(G,x)} \frac{1}{|F_n(G,x)|} d\rho = 
         \int \deg (o) h_n(G,o) d\rho,
    \end{multline}
    where $0\le h_n(G,o) \le g_{F_n}(G,o) < C$ and $\int h_n$ tends to zero in~$n$. This implies that the last integral in formula~\eqref{eq:50506} tends to zero as well. 
One can now proceed by induction in~$r$.
\end{proof}
Note that the conclusion of Proposition~\ref{proposition:rFolnler} holds if~$||g_{F_n}||_\infty$ is uniformly bounded and $||\deg||_1 < \infty$ or if $||\deg ||_\infty < \infty$ and $||g_{F_n}||_1$ is uniformly bounded (which is automatic as $||g_{A}||_1 = 1$ for any~$A$). A slightly more general sufficient condition can be deduced from the proof of Proposition~\ref{proposition:rFolnler} and the Holder inequality: it is enough to require that $||\deg||_p$ is finite and $||g_{F_n}||_q$ is uniformly bounded for some~$p, q \in [1, +\infty]$ such that $\frac{1}{p} + \frac{1}{q} = 1$. Let us also note that the identical  proof applies to the {edge-boundary}~$\partial_E F_n(G,o)$.
%~$\partial_E F_n(G,o) = \{e \in E(G) \mid e = (x,y), x \in F_n(G,o),\ y \not \in F_n(G,o)\}$ 
Similar computation to inequality~\eqref{eq:50506} shows that the condition~$\imath_E= 0$ that defines amenability (see Section~\ref{subsection_amen}), and the condition~$\imath_r =0$ are equivalent for unimodular random networks with finite expected degree and any positive integer~$r$.

%\todo[inline]{add here the argument for making $F_n$ continuous?}
Another fact that will be of use later on is that we can choose our ($r$-)F\o lner sequences to be continuous functions for every~$n$. To be precise, we prove the following proposition. 
\begin{proposition}\label{proposition:Folner_continuous}
Let~$\{\tilde F_n\}$ is a measurable sequence of subsets satisfying some of the above properties (($r$-)F\o lner, (weakly) balanced). Then there exists a measurable sequence~$\{F_n\}$ and a sequence of positive integers~$d_n$ such that for every~$(G,o) \in \mathcal{G}_*$ we have $F_n(G,o) \subset B_{d_n}(G,o)$, the subset $F_n(G,o)$ depends only on~$B_{d_n}(G,o)$, and $\{F_n\}$ satisfies the same properties as $\tilde F_n$. 
\end{proposition}

\begin{proof}
    First, we replace $\tilde F_n$ by~$F_n$ defined as 
    $$ F(G,o)= \tilde F_n(G,o) \text{ if } \tilde F_n(G,o) \subset B_{d_n}(G,o); \quad 
    F_n(G,o) = \{o\} \text{ otherwise},$$ 
    where~$d_n$ is a sufficiently large constant to be defined later. Then for any $\eps_n > 0$ and any sufficiently large~$d_n$
    \begin{equation}\label{eq:approxfolner}
        \mbbP \left( F_n(G,o) \not = \tilde F_n(G,o) \right) < \eps_n.
    \end{equation}

    We can also assume, preserving inequality~\eqref{eq:approxfolner}, that either the degrees of all vertices in $B_{d_n}$ are no more than~$d^\prime_n$, which depends on~$d_n$, or $F_n(G,o) = \{o\}$. Now, we can partition~$\mathcal{G_*}$ by the values of~$F_n(G,o)$ in~$B_{d_n}$ viewed as subsets of~$\N$ in the canonical representative of~$(G,o)$. This is a finite partition into Borel subsets. This partition can be approximated by a finite partition into the cylinder sets. This defines a function satisfying~\eqref{eq:approxfolner} and which depends only on the neighbourhood of the root of size~$R$ for some finite~$R$. We can then redefine~$d_n$ to be the maximum of~$d_n$ and the number~$R$. 

    Given relation~\eqref{eq:approxfolner} we have for any positive~$r$
    $$
     \int \frac{|\partial_r F_n (G,o)|}{|F_n(G,o)|} d\rho \le \int \frac{|\partial_r\tilde F_n (G,o)|}{|\tilde F_n(G,o)|} d\rho + \eps_n.
    $$
    Moreover~for any positive $\phi \in L^\infty(\mathcal{G}_*, \rho)$
    $$
    \int g_{F_n}(G,o)\phi(G,o) d\rho  = \int \frac{1}{|F_n(G,o)|} \sum_{x \in F_n(G,o)} \phi(G,x) d\rho.
    $$
    And, hence, 
    \begin{equation}\label{eq_contfolner}
        \left| \int g_{F_n}(G,o)\phi(G,o) d\rho  - \int g_{\tilde F_n}(G,o)\phi(G,o) d\rho \right| \le 2\eps_n ||\phi||_\infty,
    \end{equation}
    as any average of~$\phi$ does not exceed $||\phi||_\infty$. It is, therefore, enough to require~$\eps_n$ to decay to zero for preserving the property of being~$r$-F\o lner. If this decay is sufficiently fast, then the sequence~$F_n$ inherits the property of being balanced or weakly balanced. Indeed, assume~$\{\tilde F_n\}$ is balanced. One can choose $\phi (G,o)$ to be the sign of the difference $g_{F_n}(G,o) - g_{\tilde F_n}(G,o)$. Inequality~\ref{eq_contfolner} then implies convergence in $L^1(\Gr, \rho)$ of $g_{\tilde F_n}$ to~$1$. Hence, if the sequence~$\eps_n$ is summable, then we have convergence almost surely, that is, balancedness of~$\{F_n\}$. Similarly, one can treat the case of weakly balanced~$\{\tilde F_n\}$.
\end{proof}

Note that Proposition~\ref{proposition:Folner_continuous} implies, in particular, that the function $g_{F_n}$ is continuous on~$\Gr$ and depends only on $2d_n$-neighbourhood of the root.


Finally, we show that the {F\o lner} balancedness (more precisely, edge-{F\o lner} balancedness) implies the amenability of the random network in the sense of Section~\ref{subsection_amen}.

\begin{proposition}\label{proposition:balance_implies_amen}
    Suppose that~$\rho$ is an edge-{F\o lner} balanced measure on $\mathcal{G}_*$ of finite expected degree. Then~$\rho$ is amenable. 
\end{proposition}

\begin{proof}
    Due to Theorem~8.5 in~\cite{AL07} it suffices to find a sequence of measurable functions $\lambda_n \colon \Grr \to [0,1]$ such that for $\rho$~almost any~$(G,o)$
    \begin{equation}\label{eq_amen1}
        \sum_{x \in V(G)} \lambda_n(G, o, x) = 1,
    \end{equation}
    and
    \begin{equation}\label{eq_amen2}
        \lim_{n\to \infty} \int \sum_{y \sim o} \sum_{x \in V(G)} |\lambda_n(G,o,x) - \lambda_n(G,y,x)| d \rho(G,o) = 0.
    \end{equation}
    Let $\{F_n\}$ be a balanced edge-F\o lner sequence for~$\rho$. We put
    $$
    \lambda_n(G,o, x) = \frac{1}{g_{F_n}(G,o)} \cdot\frac{\I(o \in F_n(G,x))}{ |F_n(G,x)|}.
    $$
    Condition~\eqref{eq_amen1} is automatically satisfied due to the definition of $g_n$. Let us show that the second condition~\eqref{eq_amen2} holds for our choice of $\lambda_n$-s. To be precise, we will show that it is satisfied for a subsequence. First, we show that
    \begin{equation}\label{eq_amen3}
        \lim_{n\to \infty} \int \sum_{y \sim o} \sum_{x \in V(G)} |g_{F_n}(G,o)\lambda_n(G,o,x) - g_{F_n} (G,y)\lambda_n(G,y,x)| d \rho(G,o) = 0.
    \end{equation}
    Indeed, 
    \begin{multline}\label{eq_am4567567}
        \int \sum_{y\sim o}\sum_{x \in V(G)} \left| \frac{\I(o \in F_n(G,x))}{ |F_n(G,x)|} - \frac{\I(y \in F_n(G,x))}{ |F_n(G,x)|} \right| d \rho(G,o) = \\
        \int \sum_{y\sim o} \sum_{x\in V(G)} \frac{\I((o,y) \in \partial_E F_n(G,x)) + \I((y,o) \in \partial_E F_n(G,x))}{ |F_n(G,x)|} d \rho(G,o) = \\
        \int \sum_{x \in V(G)} \frac{2|\partial_E F_n(G,o)|}{|F_n(G,o)|} d \rho(G,o).
    \end{multline}
    The last equality holds due to the mass transport principle. As the sequence~$\{F_n\}$ is edge-F\o lner the result of computation~\eqref{eq_am4567567} tends to zero in~$n$. 
    Then we note that the sum 
    $$
    \sum_{y \sim o} \sum_{x \in V(G)} |\lambda_n(G,o,x) - \lambda_n(G,y,x)|
    $$
    is bounded almost surely by~$2\deg o$ as $\lambda_n(G,o, \cdot)$ is a probability measure on~$V(G)$ for every~$o \in V(G)$. Hence, it is enough to show that this sum almost surely converges to~$0$ as the degree function is integrable. We have 
    \begin{multline}
        \sum_{y \sim o} \sum_{x \in V(G)} |\lambda_n(G,o,x) - \lambda_n(G,y,x)| \le \\
        \frac{1}{g_{F_n(G,o)}}\sum_{y \sim o} \sum_{x \in V(G)} |g_{F_n}(G,o)\lambda_n(G,o,x) - g_{F_n} (G,y)\lambda_n(G,y,x)| + \\
        \frac{1}{g_{F_n(G,o)}} \sum_{y \sim o} \left|{g_{F_n}(G,y)}- {g_{F_n}(G,o)}\right| \sum_{x \in V(G)} \lambda_n(G,y,x).
    \end{multline}
    As $\{F_n\}$ is balanced, $g_n(G,o)$ tends to~$1$ almost surely. We showed that the expectation of the first sum converges to~$0$, hence, we can choose a subsequence such that is it converges to~$0$ almost surely. The sum $\sum_{x  \in V(G)} \lambda_n(G, y, x) = 1$, thus, we are left to estimate 
    \begin{multline}
        \sum_{y \sim o} \left|{g_{F_n}(G,y)}- {g_{F_n}(G,o)}\right| \le \\
        \sum_{y\sim o}\sum_{x \in V(G)} \left| \frac{\I(o \in F_n(G,x))}{ |F_n(G,x)|} - \frac{\I(y \in F_n(G,x))}{ |F_n(G,x)|} \right|.
    \end{multline}
    The expected value of this sum tends to~$0$ due to relation~\eqref{eq_amen3} and, therefore, we can again choose a subsequence to obtain a pointwise convergence to~$0$ almost surely.
\end{proof}



\section{Strong localization of eigenfunctions}

\subsection{The main results}

In this section we state and prove our main results on localization of eigenfunctions of finite-range operators as well as several corollaries of interest. 

\begin{theorem}\label{theorem:main}
    Let~$\rho \in \mathcal{U}$ be a weakly $r$-{F\o lner} balanced unimodular random network for some positive integer~$r$ and~$D$ be an operator of finite range~$r$. Suppose that there exists $\lambda\in \C$ such that~$\dim_\rho(E_\lambda^D) > 0$. Then with positive probability the network has a $\lambda$-eigenfunction with finite support.
\end{theorem}

\begin{corollary}
    If in Theorem~\ref{theorem:main} $\rho$ is extremal then the network has a finitely supported eigenfunction $\rho$-almost surely.
\end{corollary}

\begin{remark}
    One can easily show that the reverse implication of Theorem~\ref{theorem:main} always holds if the marks of the random network are taken from a discrete subspace. Suppose that for some $(G,o) \in \supp \rho \subset \mathcal{G}_*$  there exists a finitely supported $\lambda$-eigenfunction~$f$ for the operator~$D$. Then $\dim_\rho^D(E_\lambda) > 0$. Indeed, there is a positive integer~$N$ such that the support of~$f$ is contained in $B_N(G,o)$. The measure of all rooted networks that have $(N+r)$-neighbourhood isomorphic to $B_{N+r}(G,o)$, where~$r$ is the range of~$D$, is positive. Hence, the eigenspace $E_\lambda$ is non-empty with positive probability and, therefore, $\dim_\rho(E_\lambda) > 0$.   
\end{remark}

\begin{proof}[Proof of Theorem~\ref{theorem:main}]
    Let~$A(G,o) \subset V(G)$ be a measurable assignment of a finite subset in~$(G,o)$ and %We denote
   % $$
    %A^{-1}(G,o) = \{x \in G \colon o \in A(G,x)\}.
   % $$
    let~$W$ be an { invariant} subspace of~$\mathcal{H}$ of the form 
    $$
    W = \int\limits^{\oplus}_{\Gr} W{(G,o)} d \rho(G,o)  = \int\limits^{\oplus}_{\Gr} W{(G)} d \rho(G,o).
    $$
    Let $P_W = \{P_{W(G)}\} \in \Alg$ be the orthogonal projection onto~$W$. Let us define the following quantity 
    $$
    \dim_A (W) = \int\limits_{\Gr} \frac{\sum_{x \in A(G,o)} \langle P_{W(G)} \mathbb I _x , \mathbb I _x\rangle}{|A(G,o)|} d \rho(G,o).
    $$
    The mass transport principle, more precisely, formula~\eqref{eq:avfunction} applied to the function~$\langle P_{W(G)} \mathbb I _x , \mathbb I _x\rangle$, gives
    $$
    \dim_A(W) = \int g_A(G,o) \langle P_W(G)\mathbb I_o, \mathbb I_o\rangle d \rho(G,o).
    $$
    Note that here we use the fact that~$P_W(G)$ does not depend on the choice of the root in~$G$, that is, the invariance of~$W$.
    
    Let $W_A(G,o)$ be the projection of $W(G)$ to the span $\langle \mathbb I_x \colon x\in A(G,o)\rangle$, that is, the subspace of functions supported on~$A(G,o)$. One has the following inclusion  
    $$
    W{(G)} \subset  W_A(G,o) \oplus \langle \mathbb I_x \colon x\not \in A(G,o)\rangle = Y(G,o).
    $$
    Therefore, 
    \begin{equation}\label{eq_proofthm1_1}
        \frac{\sum_{x \in A(G,o)} \langle P_{W(G)} \mathbb I _x , \mathbb I _x\rangle}{|A(G,o)|} \le 
    \frac{\sum_{x \in A(G,o)} \langle P_{Y(G,o)} \mathbb I _x , \mathbb I _x\rangle}{|A(G,o)|} =
    \frac{\dim W_A (G,o)}{|A(G,o)|}.
    \end{equation}

    Now let~$W$ be the space $E_\lambda$ of $\lambda$-eigenfunctions of~$D$ and~$A$ be an element of an $r$-F\o lner sequence $\{F_n\}_n$.
    Since~$\rho$ is weakly {F\o lner} balanced by assumption, we can choose this sequence such that there exists a constant~$c > 0$ satisfying the following bound almost surely for sufficiently large~$n$
    \begin{equation}\label{eq_proof1}
        g_{F_n}(G,o) > c.
    \end{equation}
    Therefore,
    \begin{multline}\label{eq_proof2}
        \liminf_{n \to \infty} \dim_{F_n} E_\lambda = 
        \liminf_{n \to \infty} \int g_{F_n}(G,o) \langle P_{E_\lambda(G)}\mathbb I_o, \mathbb I_o\rangle d \rho(G,o) \ge \\
        c \cdot \int \langle P_{E_\lambda(G)}\mathbb I_o, \mathbb I_o\rangle d \rho(G,o) = c \Tr(P_{E_\lambda}) = c \cdot \dim_\rho (E_\lambda) > 0.
    \end{multline}
    Here, we also use that~$g_{F_n}$ is non-negative almost everywhere. On the other hand, due to inequality~\eqref{eq_proofthm1_1} we have 
    \begin{equation}\label{eq_finproj}
        \dim_{F_n} E_\lambda \le 
        \int \frac{\dim {E_{\lambda,F_n}}(G,o)}{|F_n(G,o)|} d \rho (G,o),
    \end{equation}
    where $E_{\lambda, F_n}(G,o)$ is the projection of $E_\lambda(G)$ to $\langle \mathbb I_x \colon x\in F_n(G,o)\rangle$. 

    Since~$D$ has range~$r$, the space $E_{\lambda, F_n}(G,o)$ consists of functions~$f$ supported on~$F_n(G,o)$ which satisfy $Df(x) = \lambda f(x)$ for each~$x$ at distance at least~$r$ from the boundary of~$F_n(G,o)$. Suppose that two elements~$f_1$ and~$f_2$ of~$E_{\lambda,F_n}(G,o)$ coincide on $\partial_r F_n(G,o)$. Then their difference $f_1 - f_2$ extended by~$0$ outside~$F_n(G,o)$ is a $\lambda$-eigenfunction of~$D$ with finite support in~$F_n(G,o)$. Hence, either $(G,o)$ has a $\lambda$-eigenfunction with finite support in~$F_n(G,o)$ or 
    $$
    \dim {E_{\lambda,F_n}}(G,o) \le |\partial_r F_n(G,o)|.
    $$
    Let $p_\lambda$ be the probability of the event~$B_\lambda$ that consists of $(G,o)$ that have a $\lambda$-eigenfunction with finite support. Then
    \begin{multline*}
      \int\limits_{\Gr} \frac{\dim {E_{\lambda,F_n}}(G,o)}{|F_n(G,o)|} d \rho (G,o) \le 
    \int\limits_{B_\lambda} 1 d \rho (G,o) + 
    \int\limits_{\Gr \setminus B_\lambda} \frac{|\partial_r F_n(G,o)|}{|F_n(G,o)|} d \rho (G,o) \\ \le
    p_\lambda + \int\limits_{\Gr}\frac{|\partial_r F_n(G,o)|}{|F_n(G,o)|} d \rho(G,o).  
    \end{multline*}
    Together with equations~\eqref{eq_proof2}, \eqref{eq_finproj} this gives
    \begin{equation}\label{eq_proof3}
        c \cdot \dim_\rho (E_\lambda) \le 
        \liminf_{n \to \infty}\dim_{F_n} E_\lambda \le p_\lambda + \liminf_{n \to \infty} \int\frac{|\partial_r F_n(G,o)|}{|F_n(G,o)|} d \rho(G,o).
    \end{equation}
    As the sequence~$\{F_n\}$ is $r$-F\o lner, the last summand in formula~\eqref{eq_proof3} is zero and we conclude that $p_\lambda > 0$. 
    
    In the case of extremal~$\rho$, the conclusion follows as the property of having a $\lambda$-eigenfunction is invariant under unrooted isomorphisms of the network. Hence, the probability of this event is either~$0$ or~$1$ (see~\cite{AL07}) and, as it is positive, it has to be~$1$. 
\end{proof}

    The following are the immediate corollaries of Theorem~\ref{theorem:main}. 
    
\begin{corollary}\label{corollary:algebraic}
    Let~$\rho \in \mathcal{U}$ be weakly $r$-{F\o lner} balanced and~$D$ be an  operator with finite range~$r$ and rational coefficients. Then any~$\lambda$ such that $\dim_\rho(E^D_\lambda) > 0$ is an algebraic number. 
\end{corollary}

\begin{corollary}\label{corollary:sofic1}
    Suppose that the mark space~$\Xi$ is discrete and let a sequence of networks~$\{\rho_n\}$  converge in~$\U$ to a weakly $r$-{F\o lner} balanced $\rho \in \mathcal{U}$. And let~$D$ be an operator with finite range~$r$ such that $\dim_\rho (E^D_\lambda) > 0$ for some $\lambda \in \mathbb \C$. Then for all large enough~$n$ the network $\rho_n$ has a finitely supported $\lambda$-eigenfunction for~$D$ with positive probability and, in particular, $\dim_{\rho_n}(E^D_\lambda)> 0$. 
\end{corollary}


For (strongly) {F\o lner} balanced unimodular random networks, we can prove a refined version of Theorem~\ref{theorem:main}. We will show that not only finitely supported eigenfunctions exist but that the average dimension of the subspace of $\lambda$-eigenfunctions with finite support inside a F\o lner set converges to the value of~$\dim_\rho E_\lambda$. 

\begin{theorem}\label{theorem:balanced}
    Suppose that~$\rho$ is $r$-{F\o lner} balanced and $F_n(G,o)$ is a balanced $r$-F\o lner sequence. Let~$D$ be an operator of finite range~$r$. Then 
    $$
    \dim_\rho(E_\lambda) = \lim_{n \to \infty}\int \frac{ \dim \langle f \in E_\lambda(G) \mid \supp f \subset F_n(G,o)\rangle}{|F_n(G,o)|} d\rho(G,o).
    $$
\end{theorem}

\begin{proof}
    We will follow the proof of Theorem~\ref{theorem:main} with certain changes. Instead of inequality~\eqref{eq_proof1}, we now have for $\rho$-almost every rooted network~$(G,o)$
    $$
    \lim_{n \to \infty} g_{F_n}(G,o) = 1.
    $$
    Then the inequality~\eqref{eq_proof2} becomes
    \begin{multline}\label{eq45986009457}
        \lim_{n\to \infty} \dim_{F_n} E_\lambda = 
         \lim_{n\to \infty} \int g_{F_n}(G,o)\langle P_{E_\lambda(G)}\mathbb I_o, \mathbb I_o\rangle d \rho(G,o) = \\
        \int \langle P_{E_\lambda(G)}\mathbb I_o, \mathbb I_o\rangle d \rho(G,o) = 
    \dim_\rho (E_\lambda).
    \end{multline}
    Here we use that $g_{F_n}$ is non negative for every~$n$ with~$||g_{F_n}||_{1} =1$, and the function $\langle P_{E_\lambda(G)}\mathbb I_o, \mathbb I_o\rangle$ is non-negative as well and bounded from above by~$1$. 
    Let us then estimate the dimension of $E_{\lambda,F_n}(G,o)$. Let $E_{\lambda, F_n}^*(G,o) \le E_{\lambda,F_n}(G,o)$ be the subspace of $\lambda$-eigenfunctions with finite support inside~$F_n(G,o)$. The difference of any two functions in~$E_{\lambda,F_n}(G,o)$ that coincide on $\partial_r F_n$ belongs to $E_{\lambda, F_n}^*(G,o)$. Again, we extend this difference by zero outside~$F_n(G,o)$. Therefore, 
    \begin{equation}\label{eq4589456907}
        \dim E_{\lambda,F_n}(G,o) \le \dim
     E_{\lambda, F_n}^*(G,o) + |\partial_r F_n|.
    \end{equation}
    Combining relations~\eqref{eq_finproj}, \eqref{eq45986009457} and~\eqref{eq4589456907} we obtain 
    \begin{multline}
        \dim_\rho(E_\lambda) = 
        \lim_{n\to \infty} \dim_{F_n} E_\lambda \le 
        \lim_{n\to \infty} \int \frac{\dim {E_{\lambda,F_n}}(G,o)}{|F_n(G,o)|} d \rho (G,o) \le \\
        \lim_{n\to \infty} \int \frac{ \dim  E_{\lambda, F_n}^*(G,o)}{|F_n(G,o)|} d\rho(G,o), 
    \end{multline}
    where we use the $r$-F\o lner condition 
    $$
    \lim_{n\to\infty} \int \frac{|\partial_r F_n(G,o)|}{|F_n(G,o)|} d\rho(G,o) = 0
    $$
    to obtain the last inequality.
    
    Now we show the inequality in the other direction. By definition, $  E_\lambda(G) \ge E_{\lambda, F_n}^*(G,o)$ and, hence,
    \begin{multline}
        \dim_{F_n} E_\lambda = 
        \int \frac{\sum_{x \in F_n(G,o)}\langle P_{E_\lambda(G)} \mathbb{I}_x,\mathbb{I}_x \rangle}{|F_n(G,o)|}d\rho(G,o) \ge \\
        \int \frac{\sum_{x \in F_n(G,o)}\langle P_{E^*_{\lambda, F_n}(G,o)} \mathbb{I}_x,\mathbb{I}_x \rangle}{|F_n(G,o)|}d\rho(G,o) =
        \int \frac{\dim E^*_{\lambda, F_n}(G,o) }{|F_n(G,o)|}d\rho(G,o)
    \end{multline}
    Using relation~\eqref{eq45986009457} and the fact that 
    $$
    \dim E_{\lambda, F_n}^*(G,o) = \dim
    \langle f \in E_\lambda(G) \mid \supp f \subset F_n(G,o)\rangle, 
    $$
    we obtain the reverse inequality. 
\end{proof}


\begin{thebibliography}{99}
\bibitem[]{AL07} David Aldous; Russell Lyons. Processes on Unimodular Random Networks. Electronic Journal of Probability 12 (2007) 1454 -- 1508
\bibitem[]{Ka1997amenability} Kaimanovich, Vadim. Amenability, hyperfiniteness, and isoperimetric inequalities. Comptes Rendus de l'Acad{\'e}mie des Sciences de Paris, S{\'e}rie I Math{\'e}matique 325 (1997) 999--1004
\end{thebibliography}
\end{document}
